\documentclass[11pt]{article}
\usepackage{amsmath,amssymb}
\usepackage{geometry}
\usepackage{enumitem}
\usepackage{multicol}
\geometry{margin=0.8in}
\setlength{\parindent}{0pt}
\setlength{\columnsep}{0.3in}

\begin{document}

\begin{center}
  {\Large\textbf{TMAH Precalculus Practice Packet}}\\[3pt]
  {\large Four First-Semester Sections}
\end{center}

Name: \rule{2.7in}{0.4pt}\hfill Date: \rule{1.5in}{0.4pt}

\section*{Instructions and Organization}
This packet contains original practice aligned to the first-semester sequence
in the \textit{2026--2027 Trig/Math Analysis Honors pacing calendar}.
Complete one section at a time; allow about 60 minutes per section.
Each section has a reference page and 12 problems.
\begin{enumerate}[leftmargin=*,itemsep=0.15em]
  \item Chapter 6: angles, trigonometric functions, graphs, and sinusoidal models.
  \item Chapter 7: inverse trig, equations, identities, and angle formulas.
  \item Chapter 9: polar coordinates, complex numbers, vectors, and dot products.
  \item Chapter 11: matrices, determinants, inverses, and linear systems.
\end{enumerate}
Work without a calculator. Show key steps and give exact answers, including
units in applications. An inverse-trig expression is an acceptable exact answer
when the angle is not a standard unit-circle value. Draw and label your own
axes for graphing questions; use extra paper if needed. In identity proofs,
transform the left-hand side into the right-hand side. Assume both sides are
defined. Use $k\in\mathbb Z$ for general trigonometric solutions.

\section*{Section 1 Reference: Chapter 6}
\begin{multicols}{2}
\small
\textbf{Angles and circular measure}
\[
  \theta_{\rm rad}=\theta_{\rm deg}\frac{\pi}{180},\qquad
  s=r\theta,\qquad K=\frac12r^2\theta.
\]
Use radians for arc length $s$ and sector area $K$.

\medskip
\textbf{Trigonometric functions}

For a nonzero point $(x,y)$ on the terminal ray, let $r=\sqrt{x^2+y^2}$.
\[
\begin{aligned}
  \sin\theta&=y/r,&\cos\theta&=x/r,&\tan\theta&=y/x,\\
  \csc\theta&=r/y,&\sec\theta&=r/x,&\cot\theta&=x/y.
\end{aligned}
\]
Denominators must be nonzero. On the unit circle, $r=1$.

\medskip
\textbf{Symmetry and complementary angles}
\[
\begin{aligned}
  \sin(-x)&=-\sin x,&\cos(-x)&=\cos x,\\
  \tan(-x)&=-\tan x,\\
  \sin(\pi/2-x)&=\cos x,&\cos(\pi/2-x)&=\sin x.
\end{aligned}
\]
Even: $f(-x)=f(x)$. Odd: $f(-x)=-f(x)$.

\columnbreak
\textbf{Sine and cosine graphs}

For $y=A\sin(B(x-h))+D$ or $y=A\cos(B(x-h))+D$, with $A,B\ne0$,
\[
  \text{amplitude}=|A|,\qquad T=\frac{2\pi}{|B|}.
\]
The phase shift is $h$, the midline is $y=D$, and the range is
$[D-|A|,D+|A|]$. Five key points in one period are spaced by $T/4$.

\medskip
\textbf{Other trig graphs}

Tangent and cotangent have period $\pi/|B|$; secant and cosecant have period
$2\pi/|B|$. Vertical asymptotes occur where the corresponding sine or cosine
denominator is zero.
\[
\begin{gathered}
  \sec x=1/\cos x,\qquad\csc x=1/\sin x,\\
  \cot x=\cos x/\sin x.
\end{gathered}
\]
The ranges of the untransformed secant and cosecant functions are
$(-\infty,-1]\cup[1,\infty)$.

\medskip
\textbf{Sinusoidal models}
\[
  |A|=\frac{\max-\min}{2},\qquad D=\frac{\max+\min}{2},\qquad
  \omega=\frac{2\pi}{T}.
\]
Choose a sine or cosine phase that matches the starting position.
\end{multicols}

\newpage
\begin{center}{\large\textbf{Section 1: Trigonometric Functions and Graphs}}\end{center}
\begin{enumerate}[leftmargin=*,itemsep=0pt,series=chapterSix]
  \item Convert $225^\circ$ to radians and $-7\pi/12$ to degrees.
  Find the angle in $[0,2\pi)$ coterminal with $-17\pi/6$.
  \par\vspace{\stretch{1}}

  \item A circle has radius $12$ cm. A sector has central angle $5\pi/6$.
  Find the sector's arc length and area.
  \par\vspace{\stretch{1}}

  \item The terminal ray of $\theta$ passes through $(-5,12)$.
  Find all six trigonometric function values of $\theta$.
  \par\vspace{\stretch{1}}

  \item Evaluate exactly:
  \[
    \frac{\sin(7\pi/6)\cos(5\pi/4)}{\tan(2\pi/3)}+\sec(\pi/3).
  \]
  \par\vspace*{\stretch{1}}
\end{enumerate}

\newpage
\begin{center}{\large\textbf{Section 1 --- Continued}}\end{center}
\begin{enumerate}[leftmargin=*,itemsep=0pt,resume=chapterSix]
  \item Suppose $\sin\theta=-4/5$ and $\pi<\theta<3\pi/2$.
  Find $\cos\theta$, $\tan\theta$, $\sin(-\theta)$, and
  $\cos(\pi/2-\theta)$.
  \par\vspace{\stretch{1}}

  \item For each function, state its smallest positive period and determine
  whether it is even, odd, or neither. Justify the symmetry algebraically.
  \[
    f(x)=-3\cos(2x)+1,\quad g(x)=2\sin(3x),\quad h(x)=\sin x+\cos x.
  \]
  \par\vspace{\stretch{1}}

  \item For $y=-2\sin(3x-\pi)+1$, state the amplitude, period, phase shift,
  midline, and range. Graph one full period and label five key points.
  \par\vspace{\stretch{1}}

  \item For $y=3\cos((x+\pi)/2)-2$, state the amplitude, period, phase shift,
  midline, and range. Graph one full period and label five key points.
  \par\vspace*{\stretch{1}}
\end{enumerate}

\newpage
\begin{center}{\large\textbf{Section 1 --- Continued}}\end{center}
\begin{enumerate}[leftmargin=*,itemsep=0pt,resume=chapterSix]
  \item Graph one full period of $y=\tan(2x-\pi/3)$.
  Label its zero and vertical asymptotes, and state the period.
  \par\vspace{\stretch{1}}

  \item For $y=2\sec(x/2)+1$, state the period and range. List every vertical
  asymptote in $[0,4\pi]$ and determine whether the function has any zeros.
  \par\vspace{\stretch{1}}

  \item For each function, state its period and list all vertical asymptotes
  and zeros in $[0,2\pi]$:
  \[
    y=\csc(2x),\qquad y=-\cot x.
  \]
  \par\vspace{\stretch{1}}

  \item A Ferris wheel's lowest seat height is $5$ ft and its highest is
  $65$ ft. A revolution takes $48$ seconds. A rider starts at the lowest
  point at $t=0$. Write a height model $h(t)$, find the first time the rider
  reaches the highest point, and find $h(12)$.
  \par\vspace*{\stretch{1}}
\end{enumerate}

\newpage
\begin{center}
  {\Large\textbf{Section 2: Inverse Trig, Equations, and Identities}}\\[3pt]
  Chapter 7 --- Reference and 12 Problems
\end{center}
\section*{Reference}
\begin{multicols}{2}
\small
\textbf{Inverse-trig principal ranges}
\[
\begin{aligned}
  \arcsin x&\in[-\pi/2,\pi/2],\\
  \arccos x&\in[0,\pi],\\
  \arctan x&\in(-\pi/2,\pi/2).
\end{aligned}
\]
The domains of arcsine and arccosine are $[-1,1]$; arctangent has domain
$\mathbb R$. Expressions such as $\arcsin(\sin x)$ must return an angle in
the appropriate principal range.

\medskip
\textbf{General solution patterns}
\[
\begin{aligned}
  \sin A=\sin B&:\ A=B+2\pi k\\
              &\quad\text{or }A=\pi-B+2\pi k,\\
  \cos A=\cos B&:\ A=2\pi k\pm B,\\
  \tan A=\tan B&:\ A=B+\pi k.
\end{aligned}
\]
Keep all branches and reject values outside the requested interval or original
domain. Dividing by a trig expression can lose solutions where it is zero.

\medskip
\textbf{Basic identities}
\[
\begin{aligned}
  \sin^2x+\cos^2x&=1,\\
  1+\tan^2x&=\sec^2x,\\
  1+\cot^2x&=\csc^2x.
\end{aligned}
\]
Cosine and secant are even; sine, cosecant, tangent, and cotangent are odd.
\[
  \sin(\pi/2-x)=\cos x,\qquad
  \cos(\pi/2-x)=\sin x.
\]

\columnbreak
\textbf{Sum and difference formulas}
\[
\begin{aligned}
  \sin(u\pm v)&=\sin u\cos v\pm\cos u\sin v,\\
  \cos(u\pm v)&=\cos u\cos v\mp\sin u\sin v,\\
  \tan(u\pm v)&=\frac{\tan u\pm\tan v}{1\mp\tan u\tan v}.
\end{aligned}
\]

\textbf{Double-angle formulas}
\[
\begin{aligned}
  \sin(2x)&=2\sin x\cos x,\\
  \cos(2x)&=\cos^2x-\sin^2x\\
          &=2\cos^2x-1=1-2\sin^2x,\\
  \tan(2x)&=\frac{2\tan x}{1-\tan^2x}.
\end{aligned}
\]

\textbf{Half-angle formulas}
\[
\begin{aligned}
  \sin^2(x/2)&=\frac{1-\cos x}{2},\\
  \cos^2(x/2)&=\frac{1+\cos x}{2},\\
  \tan(x/2)&=\frac{\sin x}{1+\cos x}
           =\frac{1-\cos x}{\sin x}.
\end{aligned}
\]
Square-root versions require a sign chosen from the quadrant of $x/2$.
Use identities only where their original expressions are defined.

\medskip
\textbf{Proof format}

Write a chain of equivalent expressions starting with the left-hand side.
Numerical checks alone do not prove an identity.
\end{multicols}

\newpage
\begin{center}{\large\textbf{Section 2: Practice}}\end{center}
\begin{enumerate}[leftmargin=*,itemsep=0pt,series=chapterSeven]
  \item Evaluate:
  \[
    \arcsin(-\sqrt3/2),\qquad \arccos(-\sqrt2/2),\qquad \arctan(-1).
  \]
  \par\vspace{\stretch{1}}

  \item Evaluate:
  \[
    \arcsin(\sin(5\pi/6)),\quad
    \arccos(\cos(7\pi/6)),\quad
    \arctan(\tan(3\pi/4)).
  \]
  \par\vspace{\stretch{1}}

  \item Evaluate exactly:
  \[
    \sin(\arccos(5/13)),\quad
    \tan(\arcsin(-3/5)),\quad
    \cos(2\arctan 2).
  \]
  \par\vspace{\stretch{1}}

  \item Solve on $0\le x<2\pi$:
  \[
    2\sin^2x-3\sin x+1=0.
  \]
  \par\vspace*{\stretch{1}}
\end{enumerate}

\newpage
\begin{center}{\large\textbf{Section 2 --- Continued}}\end{center}
\begin{enumerate}[leftmargin=*,itemsep=0pt,resume=chapterSeven]
  \item Solve on $0\le x<2\pi$:
  \[
    2\cos^2x+3\sin x=0.
  \]
  \par\vspace{\stretch{1}}

  \item Give all real solutions of each equation:
  \[
    \tan(2x)=\sqrt3,\qquad \cos(3x)=\cos x.
  \]
  \par\vspace{\stretch{1}}

  \item Prove:
  \[
    \frac{1-\cos(2x)}{\sin(2x)}=\tan x.
  \]
  \par\vspace{\stretch{1}}

  \item Prove:
  \[
    \sin(x+y)+\sin(x-y)=2\sin x\cos y.
  \]
  \par\vspace*{\stretch{1}}
\end{enumerate}

\newpage
\begin{center}{\large\textbf{Section 2 --- Continued}}\end{center}
\begin{enumerate}[leftmargin=*,itemsep=0pt,resume=chapterSeven]
  \item Prove:
  \[
    \frac{\cos(\pi/2+x)+\sin(-x)}{\sin(\pi/2-x)}=-2\tan x.
  \]
  \par\vspace{\stretch{1}}

  \item Suppose $\sin\alpha=3/5$ with $\alpha$ in Quadrant II, and
  $\cos\beta=-5/13$ with $\beta$ in Quadrant III.
  Find $\sin(\alpha+\beta)$ and $\cos(\alpha-\beta)$.
  \par\vspace{\stretch{1}}

  \item Suppose $\cos\theta=-7/25$ and $\pi<\theta<3\pi/2$.
  Find $\sin(2\theta)$, $\cos(2\theta)$, and $\tan(\theta/2)$.
  \par\vspace{\stretch{1}}

  \item Solve $\sin(2x)=\cos x$ on $0\le x<2\pi$.
  Which solutions would be lost by dividing both sides by $\cos x$?
  \par\vspace*{\stretch{1}}
\end{enumerate}

\newpage
\begin{center}
  {\Large\textbf{Section 3: Polar, Complex Numbers, and Vectors}}\\[3pt]
  Chapter 9 --- Reference and 12 Problems
\end{center}
\section*{Reference}
\begin{multicols}{2}
\small
\textbf{Polar coordinates and curves}
\[
  x=r\cos\theta,\quad y=r\sin\theta,\quad r^2=x^2+y^2.
\]
For polar points, signed radii are allowed:
\[
  (r,\theta)\equiv(r,\theta+2\pi k)\equiv(-r,\theta+\pi).
\]
When a positive radius is requested, choose the angle from the correct quadrant.
If a conversion divides by $r$, check the origin separately.

\medskip
For a Cartesian curve, test symmetry using $(x,-y)$ for the $x$-axis,
$(-x,y)$ for the $y$-axis, and $(-x,-y)$ for the origin. For a polar sketch,
plot key angles and check where $r=0$ and where $|r|$ is largest.

\medskip
\textbf{Complex polar form and powers}

For $z=a+bi\ne0$, the modulus is positive:
\[
  R=\sqrt{a^2+b^2},\qquad
  z=R(\cos\phi+i\sin\phi)=Re^{i\phi}.
\]
\[
  e^{i\phi}=\cos\phi+i\sin\phi,\qquad
  (Re^{i\phi})^m=R^me^{im\phi}\quad(m\in\mathbb Z).
\]
Use $0\le\phi<2\pi$ for final complex polar answers.

\medskip
\textbf{All $n$th roots of $Re^{i\phi}\ne0$}
\[
  z_k=R^{1/n}e^{i(\phi+2\pi k)/n},\qquad
  k=0,1,\ldots,n-1.
\]
The roots are equally spaced on a circle of radius $R^{1/n}$.

\columnbreak
\textbf{Vector basics}

If $P=(p_1,p_2)$ and $Q=(q_1,q_2)$,
\[
  \overrightarrow{PQ}=\langle q_1-p_1,q_2-p_2\rangle.
\]
For $\mathbf u=\langle u_1,u_2\rangle\ne\mathbf0$,
\[
  |\mathbf u|=\sqrt{u_1^2+u_2^2},\qquad
  \widehat{\mathbf u}=\frac{\mathbf u}{|\mathbf u|}.
\]
A vector of magnitude $M$ and direction angle $\theta$ has components
$\langle M\cos\theta,M\sin\theta\rangle$.

\medskip
\textbf{Dot product, angle, and projection}
\[
\begin{aligned}
  \mathbf u\cdot\mathbf v&=u_1v_1+u_2v_2,\\
  \cos\gamma&=\frac{\mathbf u\cdot\mathbf v}{|\mathbf u||\mathbf v|},\\
  \operatorname{proj}_{\mathbf u}\mathbf w
  &=\frac{\mathbf w\cdot\mathbf u}{\mathbf u\cdot\mathbf u}\mathbf u.
\end{aligned}
\]
The angle formula requires nonzero vectors. Nonzero vectors are perpendicular
exactly when their dot product is zero. The scalar component of $\mathbf F$
along nonzero $\mathbf d$ is $(\mathbf F\cdot\mathbf d)/|\mathbf d|$.

\medskip
\textbf{Applications}
\[
  W=\mathbf F\cdot\mathbf d,\qquad
  \mathbf v_{\rm ground}=\mathbf v_{\rm air}+\mathbf v_{\rm wind}.
\]
Constant-velocity displacement over time $t$ is $t\mathbf v$.
\end{multicols}

\newpage
\begin{center}{\large\textbf{Section 3: Practice}}\end{center}
\begin{enumerate}[leftmargin=*,itemsep=0pt,series=chapterNine]
  \item Convert the polar point $(-6,5\pi/6)$ to Cartesian coordinates.
  Convert $(-3,-3\sqrt3)$ to polar coordinates with $r>0$ and
  $0\le\theta<2\pi$.
  \par\vspace{\stretch{1}}

  \item Convert $r=4\cos\theta-6\sin\theta$ to Cartesian form and state
  the center and radius. Convert $(x-2)^2+y^2=4$ to a polar equation
  with $r$ isolated. Check that the origin is represented.
  \par\vspace{\stretch{1}}

  \item Convert $r^2=8\sin(2\theta)$ to Cartesian form.
  Determine whether the curve is symmetric about the $x$-axis, $y$-axis,
  and origin. Justify your decisions using the Cartesian equation.
  \par\vspace{\stretch{1}}

  \item Sketch $r=3-3\cos\theta$. Label the points for
  $\theta=0,\pi/2,\pi,3\pi/2$ and determine whether the curve is symmetric
  about the polar axis.
  \par\vspace*{\stretch{1}}
\end{enumerate}

\newpage
\begin{center}{\large\textbf{Section 3 --- Continued}}\end{center}
\begin{enumerate}[leftmargin=*,itemsep=0pt,resume=chapterNine]
  \item Write $z=-2+2\sqrt3\,i$ in trigonometric polar form and exponential
  form. State its modulus and argument in $[0,2\pi)$.
  \par\vspace{\stretch{1}}

  \item Evaluate exactly in rectangular form:
  \[
    (\sqrt3+i)^8.
  \]
  \par\vspace{\stretch{1}}

  \item Find all fourth roots of $-16$. Give every root in exponential
  and rectangular form.
  \par\vspace{\stretch{1}}

  \item Let $P=(-2,5)$, $Q=(4,-3)$, $\mathbf u=\overrightarrow{PQ}$, and
  $\mathbf v=\langle-3,4\rangle$. Find $\mathbf u$, $|\mathbf u|$, a unit
  vector in the direction of $\mathbf u$, and $2\mathbf u-3\mathbf v$.
  \par\vspace*{\stretch{1}}
\end{enumerate}

\newpage
\begin{center}{\large\textbf{Section 3 --- Continued}}\end{center}
\begin{enumerate}[leftmargin=*,itemsep=0pt,resume=chapterNine]
  \item Vector $\mathbf v$ has magnitude $12$ and direction angle $5\pi/6$.
  Vector $\mathbf w$ has magnitude $5$ and direction angle $7\pi/6$.
  Find their component forms, $\mathbf v+\mathbf w$, and $|\mathbf v+\mathbf w|$.
  \par\vspace{\stretch{1}}

  \item Let $\mathbf u=\langle2,-1\rangle$, $\mathbf v=\langle1,2\rangle$,
  and $\mathbf w=\langle3,4\rangle$. Find $\mathbf u\cdot\mathbf v$, the angle
  between $\mathbf u$ and $\mathbf v$, and
  $\operatorname{proj}_{\mathbf u}\mathbf w$.
  \par\vspace{\stretch{1}}

  \item A constant force $\mathbf F=\langle6,-2\rangle$ newtons acts through
  displacement $\mathbf d=\langle3,4\rangle$ meters. Find the work done and
  the scalar component of the force along the displacement.
  \par\vspace{\stretch{1}}

  \item Use east as the positive $x$-direction and north as the positive
  $y$-direction. A plane's air velocity is $\langle180,0\rangle$ mph and the
  wind velocity is $\langle-20,30\rangle$ mph. Find the ground velocity,
  ground speed, direction angle north of east, and displacement after
  $2$ hours. Give the direction as an exact inverse-trig expression.
  \par\vspace*{\stretch{1}}
\end{enumerate}

\newpage
\begin{center}
  {\Large\textbf{Section 4: Matrices and Linear Systems}}\\[3pt]
  Chapter 11 --- Reference and 12 Problems
\end{center}
\section*{Reference}
\begin{multicols}{2}
\small
\textbf{Matrix operations and dimensions}

Matrices can be added only when their dimensions match. If $A$ is $m\times n$
and $B$ is $n\times p$, then $AB$ is $m\times p$, with entries
\[
  (AB)_{ij}=\sum_{k=1}^{n}a_{ik}b_{kj}.
\]
In general, $AB\ne BA$. Reversing the order may even make a product undefined.

\medskip
\textbf{Determinants}
\[
  \det\begin{pmatrix}a&b\\c&d\end{pmatrix}=ad-bc.
\]
\[
\begin{aligned}
  \det\begin{pmatrix}a&b&c\\d&e&f\\g&h&i\end{pmatrix}
  &=a(ei-fh)-b(di-fg)\\
  &\quad+c(dh-eg).
\end{aligned}
\]
A square matrix is invertible exactly when its determinant is nonzero.

\medskip
\textbf{Inverses}
\[
  \begin{pmatrix}a&b\\c&d\end{pmatrix}^{-1}
  =\frac1{ad-bc}\begin{pmatrix}d&-b\\-c&a\end{pmatrix}.
\]
For a larger square matrix, row-reduce $[A\mid I]$ to $[I\mid A^{-1}]$,
if possible. Check an inverse by multiplying to obtain $I$.

\columnbreak
\textbf{Matrix equations}

A linear system can be written $A\mathbf x=\mathbf b$. If $A$ is invertible,
\[
  \mathbf x=A^{-1}\mathbf b.
\]
For $AX=B$, multiply on the left by $A^{-1}$:
\[
  X=A^{-1}B.
\]
Order matters.

\medskip
\textbf{Cramer's rule}

For a square system with $D=\det A\ne0$,
\[
  x_j=\frac{D_j}{D},
\]
where $D_j$ is the determinant obtained by replacing column $j$ of $A$ with
$\mathbf b$.

\medskip
\textbf{Row reduction and consistency}

Allowed row operations: swap two rows, multiply a row by a nonzero number,
or add a multiple of one row to another.
\begin{itemize}[leftmargin=*,itemsep=0.3em]
  \item A row $[0\ \cdots\ 0\mid c]$ with $c\ne0$ means no solution.
  \item A consistent system with a pivot for every unknown has one solution.
  \item A consistent system with a free variable has infinitely many solutions.
\end{itemize}
For a square system, a zero determinant alone does not distinguish no
solutions from infinitely many solutions.
\end{multicols}

\newpage
\begin{center}{\large\textbf{Section 4: Practice}}\end{center}
\begin{enumerate}[leftmargin=*,itemsep=0pt,series=chapterEleven]
  \item Let
  \[
    A=\begin{pmatrix}1&-2\\3&4\end{pmatrix},\qquad
    B=\begin{pmatrix}0&5\\-1&2\end{pmatrix}.
  \]
  Find $2A-B$, $AB$, and $BA$. Do $A$ and $B$ commute?
  \par\vspace{\stretch{1}}

  \item Suppose $A$ is $2\times3$, $B$ is $3\times2$, and $C$ is $2\times2$.
  For each expression, state its dimensions or explain why it is undefined:
  \[
    AB,\qquad BA,\qquad CA,\qquad AC,\qquad AB+C.
  \]
  \par\vspace{\stretch{1}}

  \item Evaluate the determinant:
  \[
    \det\begin{pmatrix}1&2&-1\\0&3&4\\2&-1&1\end{pmatrix}.
  \]
  \par\vspace{\stretch{1}}

  \item Find $A^{-1}$ and use it to solve the matrix equation:
  \[
    A=\begin{pmatrix}2&3\\1&2\end{pmatrix},\qquad
    A\begin{pmatrix}x\\y\end{pmatrix}=\begin{pmatrix}7\\4\end{pmatrix}.
  \]
  \par\vspace*{\stretch{1}}
\end{enumerate}

\newpage
\begin{center}{\large\textbf{Section 4 --- Continued}}\end{center}
\begin{enumerate}[leftmargin=*,itemsep=0pt,resume=chapterEleven]
  \item Determine all real $k$ for which
  $A=\begin{pmatrix}k&2\\3&k-1\end{pmatrix}$ is invertible.
  When $k=3$, determine the value of $b$ for which
  \[
    A\begin{pmatrix}x\\y\end{pmatrix}=\begin{pmatrix}b\\6\end{pmatrix}
  \]
  is consistent, and state how many solutions it then has.
  \par\vspace{\stretch{1}}

  \item Use Cramer's rule to solve, showing the determinants:
  \[
    2x-3y=7,\qquad 4x+y=7.
  \]
  \par\vspace{\stretch{1}}

  \item Solve using an augmented matrix and row operations:
  \[
    \begin{cases}
      x+y+z=6,\\
      2x-y+z=3,\\
      x+2y-z=2.
    \end{cases}
  \]
  \par\vspace{\stretch{1}}

  \item Classify the number of solutions for every real $t$.
  When solutions exist, give the complete solution set.
  \[
    \begin{cases}
      x+y+z=2,\\
      2x+2y+2z=t,\\
      x-y+z=0.
    \end{cases}
  \]
  \par\vspace*{\stretch{1}}
\end{enumerate}

\newpage
\begin{center}{\large\textbf{Section 4 --- Continued}}\end{center}
\begin{enumerate}[leftmargin=*,itemsep=0pt,resume=chapterEleven]
  \item Find $A^{-1}$ using an augmented matrix, then solve for $\mathbf x$:
  \[
    A=\begin{pmatrix}1&1&0\\0&1&1\\0&0&1\end{pmatrix},\qquad
    A\mathbf x=\begin{pmatrix}3\\5\\2\end{pmatrix}.
  \]
  \par\vspace{\stretch{1}}

  \item Solve the matrix equation $AX=B$:
  \[
    A=\begin{pmatrix}1&2\\3&5\end{pmatrix},\qquad
    B=\begin{pmatrix}4&1\\11&3\end{pmatrix}.
  \]
  \par\vspace{\stretch{1}}

  \item A performance sells $180$ tickets. Adult tickets cost \$12 and
  student tickets cost \$8. Total ticket revenue is \$1,824.
  Write a matrix system and solve for the number of each type of ticket sold.
  \par\vspace{\stretch{1}}

  \item Three stores sell two products. Rows of $Q$ give each store's unit
  sales; columns correspond to products 1 and 2:
  \[
    Q=\begin{pmatrix}12&8\\5&15\\10&6\end{pmatrix},\quad
    \mathbf p=\begin{pmatrix}25\\40\end{pmatrix},\quad
    \mathbf c=\begin{pmatrix}18\\28\end{pmatrix}.
  \]
  Here $\mathbf p$ and $\mathbf c$ are selling prices and unit costs in dollars.
  Find each store's revenue and gross profit using matrix products.
  Explain why $\mathbf pQ$ is not an alternative revenue calculation.
  \par\vspace*{\stretch{1}}
\end{enumerate}

\end{document}
